%%%%%%%%%%%%%%%%%%%%%%%%%%%%%%%%%%%%%%%%%
% Professional Resume - Divyanshu Goyal
%%%%%%%%%%%%%%%%%%%%%%%%%%%%%%%%%%%%%%%%%

\documentclass{resume}

\usepackage[left=0.5in,top=0.5in,right=0.5in,bottom=0.5in]{geometry}
\usepackage[T1]{fontenc}
\usepackage{xcolor}
\usepackage{hyperref}
\hypersetup{colorlinks=true, urlcolor=blue, linkcolor=blue}


\name{Divyanshu Goyal}
\address{Sunnyvale, CA \\ (+1) 404-414-3762 \\ divyanshu.g25@gmail.com}
\address{\href{https://www.linkedin.com/in/divyanshu2594/}{LinkedIn} \\ \href{https://github.com/Divyanshu25}{GitHub} \\ \href{https://scholar.google.com/citations?user=5MJm4JsAAAAJ}{Google Scholar}}

\begin{document}

%----------------------------------------------------------------------------------------
%	SUMMARY
%----------------------------------------------------------------------------------------

\begin{rSection}{Summary}
\raggedright
Applied Scientist and multimodal researcher with 10 years at the intersection of frontier ML research and production engineering. I lead the training and fine-tuning of {\bf vision-language models} at scale, design novel {\bf evaluation and alignment} methods ({\bf 5 patents}, peer-reviewed and arXiv publications), and previously led the AI powering {\bf Adobe GenStudio for Performance Marketing} (now serving Fortune 500 clients including Microsoft and AT\&T). Deep expertise across {\bf VLM training}, {\bf multimodal reasoning}, and {\bf LLM pretraining \& scaling laws}, paired with the systems engineering to deploy them in production.
\end{rSection}

%----------------------------------------------------------------------------------------
%	EXPERIENCE
%----------------------------------------------------------------------------------------

\begin{rSection}{Experience}

\begin{rSubsection}{Applied Scientist -- Adobe Unified Platform, San Jose}{June 2021 -- Present}{}{}

\item[] \vspace{-0.5em}{\textit{\textbf{$\diamond$ Vision-Language Model Training \& Evaluation}}}
\item Fine-tuned a {\bf Gemma-3-12B vision-language model} with {\bf LoRA} across {\bf 8$\times$ H100 GPUs} (DDP, bfloat16, Flash Attention 2 + Liger kernels) for large-scale marketing-quality assessment.
\item Pioneered an {\bf observable injection} technique that grounds VLM quality reasoning in {\bf 14+ computed pixel-level metrics} (sharpness, exposure, color balance); built an interpretable scorer over these signals reaching {\bf Pearson r $\approx$ 0.88} at {\bf MAE 0.77}, outperforming GPT-5.4 (MAE 1.34) and Claude Opus 4.7 (MAE 1.95) baselines.
\item Developed the {\bf MQCore evaluation benchmark} (inspired by DCLM CORE) and a fault-tolerant async annotation pipeline processing {\bf 100+ images/min} over the {\bf 700K-image KADIS} corpus.

\item[] \vspace{0.2em}{\textit{\textbf{$\diamond$ Multimodal Research \& Benchmarking}}}
\item Authored {\bf DistortBench} (CVPR FGVC13 Workshop 2026):a {\bf 13.5K-question} diagnostic benchmark spanning {\bf 27 distortion types}, 6 perceptual categories, and 5 severity levels.
\item Evaluated {\bf 18+ frontier VLMs} (GPT-5.4, Claude Sonnet 4.6, Qwen, InternVL, Gemma) with bootstrap confidence intervals and balanced-accuracy metrics, exposing low-level perception as a major weakness:top model {\bf 61.9\%} vs.\ {\bf 65.7\% human baseline}.
\item Created a curiosity-driven {\bf LLM-as-a-judge} framework for personalized creative evaluation using Bayesian surprise as a quantifiable novelty signal (arXiv 2025), achieving {\bf 23\% F1} and {\bf 38\% correlation} improvements over baselines.

\item[] \vspace{0.2em}{\textit{\textbf{$\diamond$ Production Multimodal Systems \& Agents}}}
\item Architected the {\bf Unified Brand Service} (90K+ LOC across 3 global regions) powering {\bf Adobe GenStudio}:multimodal brand-DNA extraction, RAG-based retrieval, and GPT-4-Vision content validation.
\item Built a {\bf multimodal agent} that drives Adobe Photoshop from natural language via a {\bf 49-tool} observe-think-act loop with vision grounding and an iterative quality-critic refinement layer.
\item Built the founding prototype of {\bf Adobe GenStudio for Performance Marketing} (launched Oct 2024, now serving Fortune 500 clients including Microsoft and AT\&T); invented patented VLM pipelines for {\bf cross-cultural image adaptation} and {\bf constrained machine translation} (US Patents 2026).
\end{rSubsection}

\begin{rSubsection}{Machine Learning Engineer -- Adobe Inc., Bangalore}{July 2016 -- August 2019}{}{}
\item Co-engineered Adobe's production {\bf ML inference platform}, sustaining {\bf 3,500 QPS/node} at {\bf 0.3--0.6ms} p99 latency through flame-graph profiling, JMH micro-benchmarking, and lock-free design.
\item Architected a {\bf GC-free asynchronous event pipeline} on the LMAX Disruptor ring buffer, decoupling telemetry from prediction to eliminate back-pressure at peak load.
\item Built an end-to-end {\bf model lifecycle platform}:versioned deployment, live monitoring, and staged rollout across dev, staging, and production.
\item Co-authored the \href{https://medium.com/adobetech/engineering-high-throughput-low-latency-machine-learning-services-7d45edac0271}{Adobe Tech Blog} on engineering high-throughput, low-latency ML services.
\end{rSubsection}

\end{rSection}

%----------------------------------------------------------------------------------------
%	PUBLICATIONS & PATENTS
%----------------------------------------------------------------------------------------

\begin{rSection}{Publications \& Patents}

\textbf{Curiosity-Driven LLM-as-a-judge for Personalized Creative Judgment} \hfill arXiv 2025 \\
V.\ Bannihatti Kumar, {\bf D.\ Goyal}, A.\ Eppa, N.\ Bhandari \hfill \href{https://arxiv.org/abs/2510.05135}{[paper]} \\
{\small A curiosity-driven framework that personalizes creative-writing evaluation to individual judgments via LLM-as-a-judge, targeting the subjectivity gap in LLM assessment.}

\vspace{0.6em}
\textbf{Cultural Adaptation of Images Using Generative AI} \hfill US Patent App.\ 2026 \\
A.\ Eppa, M.\ Anand, {\bf D.\ Goyal} \hfill \href{https://patents.google.com/patent/US20260051084A1/en}{[patent]} \\
{\small Multi-stage pipeline where a vision-language model grounds source images as text, conditioned on region-specific cultural guidelines, and a text-to-image model renders culturally adapted output.}

\vspace{0.6em}
\textbf{Guiding Language Translation with Translation Documents Using ML} \hfill US Patent App.\ 2026 \\
{\bf D.\ Goyal}, A.\ Eppa, M.\ Anand \hfill \href{https://patents.google.com/patent/US20260050753A1/en}{[patent]} \\
{\small Structured approach to constrained machine translation by extracting guidelines into language-agnostic facets and computing facet-specific adherence scores for linguistic and stylistic alignment.}

\vspace{0.6em}
\textbf{Parameter-Efficient Model Serving Systems} \hfill US Patent App.\ 2025 \\
Y.\ Wang, N.\ Vangala, M.\ Anand, {\bf D.\ Goyal}, et al. \hfill \href{https://patents.google.com/patent/US20250217193A1/en}{[patent]} \\
{\small System for efficiently serving fine-tuned models by dynamically loading parameter-efficient layers into pre-loaded base models, minimizing network transfer costs and container startup time.}

\vspace{0.6em}
\textbf{Brand-Aligned Content Generation with Structured Knowledge} \hfill US Patent App.\ 2025 \\
{\bf D.\ Goyal}, M.\ Anand, J.\ Mathew \hfill \href{https://patents.google.com/patent/US20250117829A1/en}{[patent]} \\
{\small Generative system that transforms unstructured brand information into structured brand DNA elements with confidence scoring, enabling brand-aligned content creation across channels.}

\end{rSection}

%----------------------------------------------------------------------------------------
%	TECHNICAL WRITING
%----------------------------------------------------------------------------------------

\begin{rSection}{Selected Projects \& Technical Writing}

\textbf{VibeNanoChat:GPT-2 Scale LLM Pretraining Framework} \hfill 2026 \\
{\small From-scratch decoder-only transformer (RoPE, RMSNorm, sliding-window attention) with one-knob depth-parameterized scaling, a custom distributed {\bf Muon} optimizer with ZeRO-2 state sharding, and Flash-Attention-3 / bfloat16 training on 8$\times$ H100. Empirically verified the {\bf C = 6ND} compute law across 40+ scaling-law runs.} \hfill \href{https://github.com/divyanshu25/VibeNanoChat}{[code]}

\vspace{0.6em}
\textbf{A Lazy Engineer's Guide to Scaling Transformers} \hfill March 2026 \\
{\small Deriving all transformer architecture dimensions, training hyperparameters, and compute schedules from a single integer (depth), enabling trivial multi-scale training sweeps.} \hfill \href{https://medium.com/@divyanshugoyal/a-lazy-engineers-guide-to-scaling-transformers-104e404e2f91}{[blog]}

\vspace{0.6em}
\textbf{Going Fast: Every Optimization That Made LLM Training Fly} \hfill March 2026 \\
{\small Six concrete optimizations (TF32, BF16 mixed precision, torch.compile, Flash Attention, parallel DataLoaders, and pinned memory) that collectively drove significant MFU improvements during GPT-2 scale training.} \hfill \href{https://medium.com/@divyanshugoyal/going-fast-every-optimization-that-made-llm-training-fly-f465f3cf3588}{[blog]}

\end{rSection}

%----------------------------------------------------------------------------------------
%	TECHNICAL SKILLS
%----------------------------------------------------------------------------------------

\begin{rSection}{Technical Skills}

\begin{tabular}{ @{} >{\bfseries}l @{\hspace{2ex}} l }
Research: & LLM Pretraining \& Alignment (SFT, RLHF, DPO), Vision-Language Models, \\
& Multimodal Agents, Evaluation \& Benchmarking, Data Curation \\
Training \& Serving: & PyTorch, Hugging Face Transformers/PEFT, DeepSpeed, FSDP, vLLM, \\
& Flash Attention, LoRA/QLoRA, CUDA, Weights \& Biases \\
Infrastructure: & Docker, LangChain, RAG pipelines \\
Programming: & Python, C++, JavaScript \\
\end{tabular}

\end{rSection}

%----------------------------------------------------------------------------------------
%	EDUCATION
%----------------------------------------------------------------------------------------

\begin{rSection}{Education}
\vspace{-0.5em}
{\bf Georgia Institute of Technology} -- M.S. Computer Science (Machine Learning) \hfill {2019 -- 2021 \textbar\ {\bf GPA: 4.0}} \\
{\bf BITS Pilani} -- B.E. Computer Science, M.Sc. Mathematics \hfill {2011 -- 2016 \textbar\ {\bf GPA: 8.75/10}}
\end{rSection}

\end{document}
